% Social-media card, 16:9.  Build:  pdflatex card.tex (twice), then
%   pdftoppm -scale-to-x 1600 -scale-to-y 900 -png -singlefile card.pdf card   ->  card.png
\documentclass{article}
\usepackage[paperwidth=32cm,paperheight=18cm,margin=0cm]{geometry}
\usepackage[T1]{fontenc}
\usepackage[utf8]{inputenc}
\usepackage[sfdefault,lining]{FiraSans}
\usepackage{amsmath}
\usepackage{newtxsf}
\usepackage{graphicx}
\usepackage{xcolor}
\usepackage{tikz}
\pagestyle{empty}
\setlength{\parindent}{0pt}

\definecolor{main}{RGB}{0,65,145}          % KAIST Dark Blue
\definecolor{accentlight}{RGB}{95,190,235} % KAIST Light Blue
\newcommand{\rr}{\raggedright\hyphenpenalty=10000 \exhyphenpenalty=10000}

% \hm{x}{y}{file}{caption}: a heatmap of side 3.55 cm with its top-left corner at (x,y)
\newcommand{\hm}[4]{%
  \node[anchor=north west,inner sep=0pt] at (#1,#2)
    {\includegraphics[width=3.55cm]{../img/#3.png}};
  \node[anchor=north,inner sep=0pt,font=\fontsize{10.5pt}{12pt}\selectfont,color=white]
    at (#1+1.775,#2-3.75) {#4};}
% \row{y}{label}: a row label
\newcommand{\row}[2]{%
  \node[anchor=south west,inner sep=0pt,font=\fontsize{12.5pt}{14pt}\selectfont\bfseries,
        color=accentlight] at (14.6,#1+0.22) {\MakeUppercase{#2}};}

\begin{document}
\begin{tikzpicture}[remember picture,overlay,shift={(current page.north west)},x=1cm,y=1cm]
  \fill[main] (0,0) rectangle (32,-18);

  % ---- left: the message
  \node[anchor=north west,inner sep=0pt] at (1.5,-1.55) {\parbox{11.6cm}{\rr
    {\fontsize{13pt}{16pt}\selectfont\bfseries\color{accentlight}\MakeUppercase{NeurIPS 2026}\par}
    \vspace{5mm}
    {\fontsize{41pt}{45pt}\selectfont\bfseries\color{white}
      Gradient descent finds extremal graphs\par}
    \vspace{7mm}
    {\fontsize{16.5pt}{22.5pt}\selectfont\color{white}
      Each square is a \textbf{graphon}, the limit of large dense graphs,
      written as a neural network and optimised directly.\par
      \vspace{4mm}
      It rediscovers classical optima, proposes candidates for open problems,
      and has led to three new theorems.\par}}};

  \node[anchor=south west,inner sep=0pt] at (1.5,-16.55) {\parbox{11.6cm}{\rr
    {\fontsize{12.5pt}{16pt}\selectfont\bfseries\color{white}
      Implicit Neural Representations for Variational Problems on Graphons\par}
    \vspace{1.5mm}
    {\fontsize{12pt}{15pt}\selectfont\color{white!80!main}
      Taeyoung Kim, Jineon Baek, Joonkyung Lee, Hongseok Yang\par}}};

  % ---- right: graphons learnt in the paper
  \row{-1.55}{Odd cycles at density $1-1/k$: proven optimal}
  \hm{14.6}{-1.55}{c5_67}{$C_5$, $p=\tfrac23$}
  \hm{18.65}{-1.55}{c5_80}{$C_5$, $p=\tfrac45$}
  \hm{22.7}{-1.55}{c7_67}{$C_7$, $p=\tfrac23$}
  \hm{26.75}{-1.55}{c7_80}{$C_7$, $p=\tfrac45$}

  \row{-7.0}{Large deviations of random graphs: new candidates}
  \hm{14.6}{-7.0}{ld_k3}{$K_3$, upper tail}
  \hm{18.65}{-7.0}{ld_c4}{$C_4$, upper tail}
  \hm{22.7}{-7.0}{ld_k13}{$K_{1,3}$, upper tail}
  \hm{26.75}{-7.0}{ld_lower}{$K_3$, lower tail}

  \row{-12.45}{New candidate structures}
  \hm{14.6}{-12.45}{h6_67}{$H_6$, $p=\tfrac23$}
  \hm{18.65}{-12.45}{h6_75}{$H_6$, $p=\tfrac34$}
  \hm{22.7}{-12.45}{h6_80}{$H_6$, $p=\tfrac45$}
  \hm{26.75}{-12.45}{pet_71}{Petersen, $p=\tfrac57$}
\end{tikzpicture}
\end{document}
